% Pipes section of the report — \input from main-vi.tex / main-en.tex.
\section{Problem 1: Pipes}

\subsection{Studying the assignment and the rules}

We studied Pipes by playing the original game at \url{https://www.puzzle-pipes.com/}
and observing its behavior. Playing first, we distilled the rules as follows:

\begin{itemize}
\item The board is an $m\times m$ grid of pipe tiles. Each tile has a \emph{fixed
shape} (straight, elbow, T-junction, cross, end) and a changeable \emph{orientation};
rotating a tile is the only move.
\item Each shape exposes ports on some of its North--East--South--West sides.
Rotation never changes the port count, only directions.
\item Goal: every open port must meet a reciprocal open port on the adjacent tile
(no port leaves the board, no dangling port), all tiles must form \textbf{one
connected} network, and the network must contain \textbf{no cycle}.
\item The water source shown in the original game is a visualization anchor: once
the network is complete, water flows from the source to every tile. Moving the
source does not change the set of valid boards.
\end{itemize}

In textbook terms this is a single-agent, deterministic, fully observable,
discrete, static task. It is a \emph{satisfaction} problem: any valid board is a
solution; minimizing manual rotations is not required (the original game does not
score rotations either).

\subsection{Search problem formulation}

Let $n=m^2$ and number the tiles in row-major order. Tile $i$ has a fixed shape
and an orientation domain $R_i$ of its distinct rotations. We formulate:

\begin{table}[h]
\centering\small
\begin{tabular}{@{}p{.20\textwidth}p{.74\textwidth}@{}}
\toprule
Component & Definition in this implementation\\
\midrule
Input & Saved scrambled board (shapes + orientations), grid size, source\\
State & Prefix $s=(a_0,\dots,a_{k-1})$ of assigned orientations\\
Initial state & Empty prefix $()$, with all tile shapes known\\
Action & Assign a boundary-valid orientation to tile $k$ (row-major order)\\
Transition & $s' = s + (a_k)$, accepted only after all feasibility checks\\
Goal test & $k=n$ and the complete board satisfies every rule\\
Cost & 1 per assignment; neither solver ranks by cost\\
Solution & One complete valid orientation tuple\\
\bottomrule
\end{tabular}
\caption{Search formulation for Pipes.}
\end{table}

One subtlety: a \emph{solver step} assigns an absolute orientation to the next
tile, while a \emph{player step} rotates a tile by 90 degrees. One assignment can
equal zero to three clicks. Both solvers share the assignment model, which makes
their comparison fair. The saved scrambled board only supplies the domains $R_i$
and the rotation order; search always starts from the empty prefix and never
reads a stored solution.

\subsection{Search space and shared pruning}

The number of complete assignments before boundary checks is
$U=\prod_i |R_i| \le 4^n$: elbows, T-junctions and ends have 4 distinct rotations,
straights 2, crosses 1. A typical 4$\times$4 board has
$U = 4^{16} \approx 4.29\times10^9$; the 7$\times$7 upper bound is
$4^{49} \approx 3.17\times10^{29}$. The assignment tree has depth $n$, branching
factor at most 4, and \textbf{every prefix has exactly one parent} (fixed
assignment order), so no closed set is needed for duplicate detection.

Orientations pointing off the board are removed up front. For each candidate
child, the shared \code{compatible()} routine applies three sound pruning layers
to both solvers:

\begin{enumerate}
\item \textbf{Match assigned neighbors}: every unassigned domain $R_i$ is
restricted to orientations matching the assigned neighbors' ports (present and
absent alike). An empty domain rejects the branch.
\item \textbf{Cycle rejection}: if assigned tiles already form a cycle, reject
(disjoint-set check).
\item \textbf{Possible connectivity}: build an \emph{optimistic} graph over tiles
--- an edge exists if some remaining orientation could open toward the other
tile. If this graph cannot connect every tile, reject.
\end{enumerate}

At the root we check the port count: a tree over $n$ vertices has exactly $n-1$
edges $=2(n-1)$ ports, and rotation preserves the count, so a different total
proves unsolvability immediately (the converse does not hold). All three layers
only remove states that cannot lead to a solution, so completeness is preserved.
DFS is therefore uninformed in its \emph{traversal order} but is a
constraint-pruned search, not a blind rotation enumerator.

\subsection{Blind search: Depth-First Search}

DFS follows the first feasible child until it reaches a goal or a dead end, then
backtracks. Rotations are enumerated in each tile's saved-orientation order ---
no heuristic information is used:

\begin{lstlisting}
SEARCH_DFS(input):
  build boundary-valid domains R_i
  if root feasibility fails: return UNSOLVABLE
  return VISIT(())

VISIT(s):
  if budget exhausted: return LIMIT_REACHED
  count one expansion
  if GOAL(s): return SOLVED(s)
  for a in boundary-valid rotations of tile |s|:
    child = s + (a)
    if FEASIBLE(child):           # 3 pruning layers
      r = VISIT(child)
      if r is SOLVED or LIMIT_REACHED: return r
  return UNSOLVABLE
\end{lstlisting}

The prefix space is finite and acyclic and the pruning is sound, so DFS is
\emph{complete} given enough resources. The first solution found does not
minimize assignments. Worst-case time is exponential in $n$; memory keeps the
current path plus pending alternatives --- this Python implementation stores
copied prefix tuples and generator frames, about $O(n^2)$ along one branch.

\subsection{Heuristic search: Greedy Best-First}

Greedy always pops from a min-heap the state with the fewest \emph{remaining
choices}:
\[
h(s) = \sum_{i \text{ unassigned}} \log_2 |D_i(s)|, \qquad
\text{priority}(s) = (h(s),\ -\text{depth}(s),\ \text{insertion order})
\]
where $D_i(s)$ is the boundary-valid domain of tile $i$ after matching assigned
neighbors. The sum is the logarithm of the product of remaining domain sizes:
domains $(1,2,4)$ give $h=0+1+2=3$, while $(1,1,2)$ give $h=1$ --- state B is
less ``free'' and is preferred. We chose this heuristic because (i) it is cheap,
reusing domains already computed during pruning, and (ii) every solution at the
same depth needs exactly $n$ assignments, so $n-k$ distinguishes nothing, while
the log-domain sum prioritizes branches that narrow fastest. $h$ is
\emph{not admissible} --- it measures remaining uncertainty, not a cost bound ---
and no $g$ term is used.

\begin{lstlisting}
SEARCH_GREEDY(input):
  build boundary-valid domains; root feasibility check
  push () into min-priority queue with priority 0
  while queue is not empty:
    if budget exhausted: return LIMIT_REACHED
    s = pop minimum-priority state;  count one expansion
    if GOAL(s): return SOLVED(s)
    for a in boundary-valid rotations of tile |s|:
      child = s + (a);  D = FEASIBLE(child)
      if D exists: push child with (h(child), -|child|, serial)
  return UNSOLVABLE
\end{lstlisting}

On the finite, acyclic prefix space, keeping all feasible candidates makes Greedy
complete given enough resources; but Greedy must hold its frontier in the heap,
so memory can grow exponentially.

\subsection{Implementation and visual demo}

The implementation lives in the \code{pipes\_problem} package: \code{model.py}
(rules, validation, board generation), \code{search.py} (both solvers, pruning,
timing and memory), \code{app.py} (pygame interface), \code{worker.py} (isolated
measurement process), \code{benchmark.py} (batch runner). The demo interface
offers:

\begin{itemize}
\item Manual play: click to rotate (Shift for counterclockwise), right-click to
lock, Undo/Reset; connected pipes carry animated blue water from the source.
\item \textbf{Step-by-step search demo}: pick DFS or Greedy, then \emph{Run /
Pause / Next / Speed}. Each event is labeled: \emph{expand} (DFS visits / Greedy
pops a prefix), \emph{assign} (DFS chooses a child), \emph{queue} (Greedy pushes),
\emph{reject} (a pruned orientation --- tile outlined in red with the reason),
\emph{backtrack}. \emph{Next} separates expansion from successor generation.
\item \emph{Show solution} replays the final assignment on its own; the board
returns to its saved state before a run. \emph{Unsolvable} and \emph{limit
reached} are distinct, explicit messages.
\end{itemize}

\subsection{Creating inputs}

Inputs are produced by the project generator (not transcribed from the web):
nine default maps covering sizes $4\times4, 5\times5, 7\times7$ with seeds
261--263, stored as compact JSON (\code{size}, \code{source}, \code{tiles}) with
the filename as the input ID. In addition, \code{maps/custom/cdai.json} is a
hand-edited \emph{unsolvable} map used to check messaging and exhausted-space
behavior. Runtime generation by \code{-{}-size}/\code{-{}-seed} is also
available. These are the team's own examples (not official website puzzle IDs),
and the generator never stores a solution in the map file.

\subsection{Experimental results: time and memory}

\textbf{Protocol.} Runs were executed on Linux x86\_64 (kernel 6.8), Python
3.13.14, pygame-ce 2.5.8. Both algorithms receive the \emph{same} nine maps, the
same row-major tile order, the same pruning, a 30-second limit and 200{,}000
expansions per run. Each (map, algorithm) pair runs 5 times: time is a headless
\code{perf\_counter} solve (no drawing); memory is the peak Python allocation
measured by \code{tracemalloc} in a separate replay, combined only when both
passes agree on status, returned board and expansion count. The full 90-run
dataset is stored at \code{docs/benchmarks/2026-10-25-pipes.csv}; the tables
below report medians per input.

\begin{table}[h]
\centering\small
\begin{tabular}{@{}llrrr@{}}
\toprule
Input & Algorithm & Solved/runs & Time (ms) & Memory (KiB)\\
\midrule
pipes-4x4-261 & DFS    & 5/5 & 3.380  & 26.836\\
              & Greedy & 5/5 & 2.866  & 11.695\\
\midrule
pipes-4x4-262 & DFS    & 5/5 & 3.047  & 27.078\\
              & Greedy & 5/5 & 3.386  & 12.359\\
\midrule
pipes-4x4-263 & DFS    & 5/5 & 3.149  & 26.984\\
              & Greedy & 5/5 & 2.741  & 12.500\\
\midrule
pipes-5x5-261 & DFS    & 5/5 & 5.352  & 53.203\\
              & Greedy & 5/5 & 7.094  & 19.500\\
\midrule
pipes-5x5-262 & DFS    & 5/5 & 5.906  & 53.344\\
              & Greedy & 5/5 & 5.882  & 19.219\\
\midrule
pipes-5x5-263 & DFS    & 5/5 & 4.564  & 53.297\\
              & Greedy & 5/5 & 5.712  & 19.219\\
\midrule
pipes-7x7-261 & DFS    & 5/5 & 23.247 & 145.852\\
              & Greedy & 5/5 & 27.049 & 27.594\\
\midrule
pipes-7x7-262 & DFS    & 5/5 & 21.938 & 145.930\\
              & Greedy & 5/5 & 26.669 & 27.070\\
\midrule
pipes-7x7-263 & DFS    & 5/5 & 40.785 & 145.852\\
              & Greedy & 5/5 & 38.246 & 29.094\\
\bottomrule
\end{tabular}
\caption{Pipes --- time and peak memory per input (median of 5 runs, measured
25 October 2026). All 90/90 runs are \emph{solved}. Source:
\code{docs/benchmarks/2026-10-25-pipes.csv}.}
\label{tab:pipes-per-input-en}
\end{table}

\begin{table}[h]
\centering\small
\begin{tabular}{@{}llrrr@{}}
\toprule
Size & Algorithm & Solved/runs & Median time (ms) & Median memory (KiB)\\
\midrule
$4\times4$ & DFS    & 15/15 & 3.149  & 26.984\\
           & Greedy & 15/15 & 3.097  & 12.359\\
$5\times5$ & DFS    & 15/15 & 5.261  & 53.297\\
           & Greedy & 15/15 & 6.157  & 19.219\\
$7\times7$ & DFS    & 15/15 & 24.469 & 145.852\\
           & Greedy & 15/15 & 31.695 & 27.594\\
\bottomrule
\end{tabular}
\caption{Pipes --- summary by board size (3 maps $\times$ 5 runs per figure).}
\label{tab:pipes-summary-en}
\end{table}

\textbf{Interpreting the results.}

\begin{itemize}
\item \textbf{Time grows quickly with size} (3--5 ms at 4$\times$4 to 22--41 ms
at 7$\times$7): variables grow from 16 to 49 and each successor's feasibility
check touches more assigned neighbors, matching the exponential prediction for
the search space.
\item \textbf{DFS is faster than Greedy at 5$\times$5 and 2 of 3 maps at
7$\times$7.} On easy inputs the pruning already steers DFS's first branch almost
straight to a solution; the priority queue only adds heap push/pop and $h$
computation without changing traversal much. This shows that a heuristic is
\emph{not free}: its bookkeeping cost must be repaid by ranking quality.
\item \textbf{On the hardest map (pipes-7x7-263) Greedy wins} (38.2 vs 40.8 ms):
this map forces real backtracking, and ranking by ``fewest remaining choices''
explores tight branches first, saving some wasted probing. The margin is small
(5\%) but appears exactly where heuristics should help.
\item \textbf{Memory: Greedy uses 4--5$\times$ less than DFS} (12--29 KiB vs
27--146 KiB). Counter to the textbook worry about Greedy's exponential frontier,
the tracemalloc figure reflects this implementation's actual Python allocations:
DFS copies a prefix tuple and holds a generator frame per recursion level, while
Greedy's frontier stays tiny on easy inputs because pruning discards almost
everything. On much harder inputs Greedy's frontier could still balloon ---
further experiments would be needed to confirm.
\item \textbf{All 90 runs solved} within limits; no run hit the 30-second or
200k-expansion budget on this input set.
\end{itemize}

\subsection{Limitations}

The default input set is quite easy under the current pruning, so neither solver
is pushed to its limits; any ``winner'' claim holds only for this kind of input.
We have not yet ablated the pruning layers separately, added many unsolvable
maps, or tried an MRV heuristic. These limitations are recorded in
\code{docs/pipes-analysis.md} as future work.
